\chapter{Introducción}

\section{Motivación}

El desarrollo de robots bípedos ha ganado gran interés en las últimas décadas debido a su potencial para operar en entornos diseñados para humanos. A diferencia de los robots con ruedas, los robots bípedos pueden navegar por escaleras, terrenos irregulares y espacios estrechos. Este proyecto se centra en el control de una pierna robótica bípeda de 3 grados de libertad (DOF) utilizando controladores ODrive, sentando las bases para el desarrollo de un robot humanoide completo.

\section{Objetivos del Proyecto}

Los objetivos principales de este Trabajo de Fin de Grado (TFG) son:

\begin{itemize}
    \item Diseñar e implementar el control de una pierna robótica de 3-DOF utilizando controladores ODrive v3.6
    \item Desarrollar el software de control en Python integrado con ROS2 (distribución Jazzy)
    \item Implementar cinemática directa e inversa para la pierna robótica
    \item Generar trayectorias de marcha para el movimiento del robot
    \item Realizar pruebas experimentales de los motores y verificar el funcionamiento del sistema
\end{itemize}

\section{Estructura de la Memoria}

Esta memoria se estructura en los siguientes capítulos:

\begin{enumerate}
    \item \textbf{Introducción}: Motivación, objetivos y estructura del documento
    \item \textbf{Marco Teórico}: Conceptos fundamentales de motores BLDC, control ODrive, cinemática y generación de trayectorias
    \item \textbf{Diseño del Sistema}: Selección de hardware, arquitectura del sistema y dimensionado
    \item \textbf{Implementación}: Desarrollo del software, comunicación con ODrive y cinemática
    \item \textbf{Resultados}: Pruebas realizadas y verificación del sistema
    \item \textbf{Conclusiones y Trabajo Futuro}: Conclusiones obtenidas y líneas de trabajo futuras
    \item \textbf{Anexos}: Datasheets, código fuente y diagramas
\end{enumerate}
